\chapter{Primitive disorder, the scalar process beyond the support, and the parameters}\label{u:sec:setup}

Throughout \cref{part:upper}, we fix $\beta>1$. This section sets
the notation for that Part. We identify the
primitive Gaussian disorder with the completed model of \cref{sec:model},
record the Gaussian matrix estimates used in
\cref{u:sec:proposal,u:sec:connections,u:sec:proof}, and give the exponent
parameters. We then continue the scalar diffusion of \cref{sec:inputs}
beyond the support of the Parisi measure. The scalar profiles and the
one-row law defined below are used throughout the construction of the
proposal and its analysis. We follow the scalar construction of
\cite[Sections~2--3]{oai-typical-upper}.

The notation in \cref{u:tab:local-notation} is local to this Part.
All other notation retains the meaning fixed in \cref{sec:model,sec:inputs}.
In particular, $\PP$ denotes probability, $P_\beta$ is the limiting
pressure, and all exponents are written in sans serif.

\begin{table}[htbp]
\centering
\small
\caption{Notation local to \cref{part:upper}. The right column records
the meanings used here and any earlier meanings that they replace.}
\label{u:tab:local-notation}
\begin{tabular}{@{}>{\raggedright\arraybackslash}p{0.30\textwidth}>{\raggedright\arraybackslash}p{0.66\textwidth}@{}}
\toprule
Symbols & Meaning in this Part\\
\midrule
$U,m$ & Scalar potential and its first derivative in variance time;
these replace $\Phi_{\alpha_\beta},u_{\alpha_\beta}$ after the time change.\\
$Y,M$ & Scalar diffusion and its Wiener driver in variance time;
$Y$ replaces the overlap-time diffusion of \cref{sec:inputs}.\\
$B,D$ & Reflected Wiener pair; $B$ replaces the driver denoted $B$
in overlap time.\\
$\delta,\epsilon_n$ & Smoothing variance $n^{-\ex D}$ and cut-balance
threshold $n^{-\ex d}$.\\
$h,h_j$ & Mesh lengths; fields involving the completed matrix are
written $(W(g)x)_i$.\\
$\rho,\sigma$ & Deterministic reflection coefficients, replacing the
overlap scale $\rho$ and the resampled spin symbol $\sigma$.\\
$\theta$ & Rotation angle, or an exponential-moment parameter in the
separate concentration arguments.\\
$p,r,e,j$ & Scalar profiles, with $e_h,b_j,s_j,w_j$ their mesh quantities;
the subscripted probabilities $p_i^J$ keep their meaning.\\
$R_V,R_h,R_n,R_D$ & Inverse-profile bound, mesh residual, enlargement
radius, and reference law, respectively; $R(x,y)$ remains the overlap.\\
$H,H_0,E_H$ & Query cross-moment matrix, its population value, and
its error; Hamiltonians carry their superscripts or subscripts.\\
$H_*$ & Maximum absolute spin energy in \cref{u:sec:fine}, and a
Bernstein scale in \cref{u:sec:concentration}.\\
$\mu^{\rm src}_0,\mu^{\rm rep}_0$ & Unlinked source and replay laws.\\
$\Lambda^{\rm src},\Lambda^{\rm rep}$ & Linked source and replay laws.\\
$Q_V,Q_F,\nu_V$ & Output law, fine output law, and one-row law.\\
$L$ & Spin likelihood in \cref{u:sec:proposal,u:sec:fine};
$L=\Upt(E_A)$ in \cref{u:sec:concentration,u:sec:coarse}.\\
$W_{\rm wt},W_0,W_2$ & Weighted error size, ambient GOE matrix, and
$2$-Wasserstein distance; $W(g)$ is the completed matrix.\\
$\mathcal C_n,\mathcal H_n,\mathcal W_n,\mathrm{Bad}$ & Events and spin sets defined
where they are used; $\mathcal H_n$ denotes the pressure event.\\
\bottomrule
\end{tabular}
\end{table}

\section{Primitive disorder and the completed Hamiltonian}\label{u:sec:primitive}

Let $\gamma_n$ be standard Gaussian measure on $\R^{n\times n}$.
For $g\in\R^{n\times n}$ and $x\in\Sigma_n$, define
\begin{equation}\label{u:eq:primitive}
\begin{gathered}
 W(g)=\frac{g+g^{\mathsf T}}{\sqrt{2n}},\qquad
 H_g^c(x)=\frac12x^{\mathsf T}W(g)x,\\
 F^c(g)=\log\sum_{x\in\Sigma_n}e^{\beta H_g^c(x)},\qquad
 \pi_g(x)=e^{\beta H_g^c(x)-F^c(g)}.
\end{gathered}
\end{equation}
The dictionary to \cref{sec:model-completion} is
\[
 J_{ij}(g)=\frac{g_{ij}+g_{ji}}{\sqrt2}\quad(i<j),\qquad
 J'_{ii}(g)=g_{ii},\qquad A_{ij}(g)=\frac{g_{ij}-g_{ji}}2\quad(i<j).
\]
We note that the variables $J_{ij}$ and $J'_{ii}$ are independent
standard Gaussians.
The antisymmetric coordinates $A_{ij}$ are independent $N(0,1/2)$
variables, independent of $W(g)$. Expanding the quadratic form, we find
\[
 H_g^c(x)=H_n^{J(g)}(x)+G(g),\qquad
 G(g)=\frac12\tr W(g)=\frac1{\sqrt{2n}}\sum_i g_{ii}\sim N(0,1/2).
\]
Thus $F^c(g)=F_n^c$, and \cref{lem:model-completion}(b) implies that
$\pi_g=\pi^{J(g)}$. We write $\mathcal G_n=W_*\gamma_n$ for the GOE
law with off-diagonal variance $1/n$ and diagonal variance $2/n$, and
$\Unif_n(x)=2^{-n}$ for the uniform law on $\Sigma_n$.

For the dynamics, let $P_t^g=P_t^{J(g)}$. The rate of a spin flip is
\begin{equation}\label{u:eq:flip-rate}
 p_i^{J(g)}(-x_i\mid x)
 =\frac{\pi_g(x^{i,-x_i})}{\pi_g(x)+\pi_g(x^{i,-x_i})}.
\end{equation}
The generator is $\mathcal L^{J(g)}$, with rate one per site as in
\eqref{eq:common-semigroup}. For a cut $S\subseteq\Sigma_n$, its
Dirichlet flow is
\begin{equation}\label{u:eq:cut-flow}
 \mathcal Q_g(S,S^c)=\sum_{x\in S}\sum_{i=1}^n
 \pi_g(x)p_i^{J(g)}(-x_i\mid x)
 \mathbf 1\{x^{i,-x_i}\notin S\}.
\end{equation}
The diagonal of $W(g)$ adds the same constant to all spin energies,
so it cancels from the Gibbs measure and the dynamics.
We compare proposals with the joint law
\begin{equation}\label{u:eq:target-law}
 \Pi_n(\dd g,x)=\gamma_n(\dd g)\pi_g(x).
\end{equation}
Its $(J(g),x)$-marginal is $\PP_{J,X_0\sim\pi^J}$ from
\cref{def:model-joint}. In the transition law, the spin sampled from
$\pi_g$ is held fixed.

The relative entropy is
$\KL(\mu\Vert\nu)=\int\log(\dd\mu/\dd\nu)\dd\mu$ when
$\mu\ll\nu$, where the integral takes values in the extended
nonnegative reals, and $+\infty$ otherwise. In finite sums, $0\log0=0$, and a positive numerator over
zero contributes $+\infty$. We use the asymptotic conventions of
\cref{sec:model-asymptotics}. In every eventual bound, we first fix
$\beta>1$ and all exponents, then choose finite constants and $n_0$,
and then assert the bound for every integer $n\ge n_0$. Constants $C_\beta,c_\beta>0$ may change between
inequalities. They depend only on $\beta$ and, when an exponent is
absorbed, on the displayed fixed exponents; they are independent of
$n$, the selected cut, and the rotation angle. The bounds are for
each fixed temperature and saving $\omega<1/4$.

\paragraph{Conventions.}\label{u:par:conventions}
Spaces of continuous paths carry the Borel sigma algebra of the
uniform topology, products carry the product sigma algebras, and
Wiener measures are supported on paths starting at zero. A bridge of duration $h>0$ has the law $\mathrm{Br}_h$
obtained from Wiener measure by $b(t)\mapsto b(t)-(t/h)b(h)$.
These spaces are standard Borel. We complete filtrations when
constructing stochastic integrals and then choose Borel versions on
the original spaces. All refinements defining integrals are taken at
fixed $n$, fixed horizon, and fixed query mesh, before $n\to\infty$.
Versions are identified outside fixed-$n$ null sets. For the terminal
spin likelihood $L$ defined in \eqref{u:eq:spin-link}, the bound
$L\le2^n$ is used only to transfer fixed-$n$ null sets and integral
versions. All probability estimates are proved under the law to which
they apply.

\section{Gaussian orthogonal ensemble facts}\label{u:sec:goe}

We record the Gaussian facts needed for the adaptive construction of
the matrix in \cref{u:lem:assembly} and for the exceptional events of
the disorder in \cref{u:sec:connections,u:sec:proof}. The successive
conditioning on orthogonal queries is the Gaussian-revelation method
used in \cite[Sections~2--3]{bolthausen-2014} and in the proposal of
\cite[Section~4.1]{oai-typical-upper}.

\begin{lemma}[Gaussian conditioning]\label{u:lem:goe-conditioning}
Let $W\sim\mathcal G_n$. Its law is invariant under orthogonal
conjugation. In an orthonormal basis, the entries on and above the
diagonal are independent, with variances $2/n$ on the diagonal and
$1/n$ off the diagonal. After revealing its columns in a fixed
subspace, the compression to the orthogonal complement is an
independent GOE block with the original variance normalization $1/n$.
The same conditional assertion holds for successive orthogonal unit
query directions chosen measurably from the previously revealed columns
and independent initial data.
\end{lemma}

\begin{proof}
For deterministic vectors $u,v,u',v'$, the covariance is
\[
 \E[(u^{\mathsf T}Wv)(u'^{\mathsf T}Wv')]
 =\frac1n\bigl[(u\cdot u')(v\cdot v')
                 +(u\cdot v')(v\cdot u')\bigr].
\]
This follows by expanding \eqref{u:eq:primitive}. Since $W$ is
Gaussian, this covariance implies that its law is invariant under
orthogonal conjugation and that its entries on and above the diagonal
are independent in any fixed orthonormal basis. In particular, the unrevealed block is independent
of the columns already read. Conditional on the past, the next query
direction is fixed in the remaining subspace. In a basis starting
with this direction, its diagonal entry is $N(0,2/n)$, its remaining
column entries are independent $N(0,1/n)$, and the smaller block is
independent of both. By induction, the assertion holds for
adaptive queries, including conditioning on independent initial data.
\end{proof}

\begin{lemma}[Operator norm]\label{u:lem:goe-norm}
For every $n\ge2$,
\begin{equation}\label{u:eq:goe-norm}
 \gamma_n\{g:\norm{W(g)}_{\mathrm{op}}>12\}\le2e^{-n}.
\end{equation}
\end{lemma}

\begin{proof}
We use the sphere-net argument of
\cite[Corollary~4.2.13 and Exercise~4.4.3(b)]{vershynin-2018}.
Let $\mathcal N$ be a $1/4$ net of the Euclidean unit sphere with
$|\mathcal N|\le9^n$, obtained by comparing the volumes of disjoint
balls. For a symmetric matrix, approximating a unit vector at which
the absolute value of the quadratic form is maximal by a point of the
net, we have
$\norm W_{\mathrm{op}}\le2\max_{u\in\mathcal N}|u^{\mathsf T}Wu|$.
For every fixed unit $u$, the quadratic form $u^{\mathsf T}Wu$ is
centered Gaussian with variance $2/n$. Using the Gaussian tail and a
union bound, we find that
\[
 \gamma_n\{\norm W_{\mathrm{op}}>12\}
 \le2\,9^n e^{-9n}
 =2e^{(\log9-9)n}\le2e^{-n}.
\]
\end{proof}

\begin{lemma}[Pressure tail]\label{u:lem:pressure-tail}
Fix $\beta>1$ and $0<\ex z_{\rm P}<1/2$. There exists $c_\beta>0$
such that
\begin{equation}\label{u:eq:pressure-tail}
 \gamma_n\{F^c(g)>nP_\beta+n^{1-\ex z_{\rm P}}\}
 \le e^{-c_\beta n^{1-2\ex z_{\rm P}}}=o(1).
\end{equation}
\end{lemma}

\begin{proof}
Differentiating with respect to the primitive entries, we have
\[
 \partial_{g_{ij}}F^c(g)
 =\frac\beta{\sqrt{2n}}\sum_x\pi_g(x)x_ix_j,\qquad
 \norm{\nabla F^c(g)}^2\le\frac{\beta^2n}{2}.
\]
By \cref{inp:parisi} and \cref{lem:model-completion}(c),
$\E_{\gamma_n}F^c=\E F_n\le nP_\beta$. Then \eqref{u:eq:pressure-tail}
follows from Gaussian concentration
\cite{tsirelson-ibragimov-sudakov-1976}, with the Lipschitz constant
given by the gradient bound.
\end{proof}

\section{The exponent parameters}\label{u:sec:parameters}

Fix $\omega\in(0,1/4)$ and define
\begin{equation}\label{u:eq:epsilon-omega}
 \varepsilon_\omega=\frac{1/4-\omega}{20}\in(0,1/80).
\end{equation}
We use the exponents in \eqref{u:eq:exponents} throughout this Part.
The strict inequalities between them are verified in
\cref{u:sec:budget}, where
we also give the instance $\omega=1/5$. The letters $\ex r$ and
$\ex w$ denote exponents, whereas $r(t)$ is a scalar profile and
$\gamma_n$ is Gaussian measure.

\begin{equation}\label{u:eq:exponents}
\begin{array}{@{}cc@{\qquad}cc@{}}
\toprule
\text{Exponent} & \text{Value} & \text{Exponent} & \text{Value}\\
\midrule
\ex a & 1/2 & \ex b & 1/2-14\varepsilon_\omega\\
\ex f & 1/2 & \ex z & 1/4-8\varepsilon_\omega\\
\ex k_{\rm f} & 1/2-2\varepsilon_\omega & \ex w & 1/8-3\varepsilon_\omega\\
\ex D & 1/4-\varepsilon_\omega & \ex k_{\rm b} & 1/4-9\varepsilon_\omega\\
\ex u_{\rm E} & 1/4-2\varepsilon_\omega & \ex s_{\rm b} & 1/4-5\varepsilon_\omega\\
\ex d & \varepsilon_\omega/2 & \ex s_{\rm seq} & 1/4-4\varepsilon_\omega\\
\ex e & 1/4-3\varepsilon_\omega & \ex r & 1/4-15\varepsilon_\omega/2\\
\ex u_{\rm F} & 1/4-39\varepsilon_\omega/4 & \ex q & 1/4-10\varepsilon_\omega\\
\ex z_{\rm P} & 3/10 & \ex g_{\rm rot} & 3/5\\
\bottomrule
\end{array}
\end{equation}

\section{The scalar Parisi process and its continuation}\label{u:sec:scalar}

We use \cref{inp:parisi,inp:unique,inp:full-rsb,inp:ob-scalar,cor:inp-identity}
and the scalar evolution identities of \cref{sec:scalar}.
Let $\mu_\beta$ be the unique minimizing Parisi measure and let
$q_\beta\in(0,1)$ be its maximal support point. In variance time
$v=\beta^2t$, define
\[
 v_*=\beta^2q_\beta,\qquad
 \bar\alpha(v)=\alpha_\beta(v/\beta^2),\qquad
 U(v,y)=\Phi_{\alpha_\beta}(v/\beta^2,y)
 \quad(0\le v\le\beta^2).
\]
Then $\partial_vU+U_{yy}/2+\bar\alpha U_y^2/2=0$, with terminal
value $\log(2\cosh y)$ at $\beta^2$. We continue the potential and the
distribution function by
\begin{equation}\label{u:eq:scalar-continuation}
 U(v,y)=\log(2\cosh y)+(\beta^2-v)/2,\qquad
 \bar\alpha(v)=1\quad(v\ge v_*).
\end{equation}
This agrees with the original solution on $[v_*,\beta^2]$ and
satisfies the same equation for all later times. Let $m=U_y$ and
$\chi=U_{yy}$. Up to $\beta^2$, these are
$u_{\alpha_\beta}(v/\beta^2,y)$ and
$\chi_{\alpha_\beta}(v/\beta^2,y)$, and for $v\ge v_*$ they are
$\tanh y$ and $\operatorname{sech}^2y$. We write
$b(v,y)=\bar\alpha(v)m(v,y)$ before $v_*$ and $b=m$ afterward.
The value of the distribution function at a time atom does not
affect its time integrals.

On canonical Wiener space with driver $M$, let
\begin{equation}\label{u:eq:scalar-process}
 Y_t=M_t+\int_0^t b(v,Y_v)\dd v,\qquad
 m_t=m(t,Y_t),\qquad \chi_t=\chi(t,Y_t).
\end{equation}
Here and below, the time argument of $Y$ is variance time. Up to
$\beta^2$, the dictionary to \eqref{b:eq:scalar-diffusion} is
$Y^{\rm IV}_v=Y_{v/\beta^2}$ and $M_v=\beta B_{v/\beta^2}$.
Since the drift is bounded and spatially $1$-Lipschitz, there exists
a unique continuous, measurable, causal Picard solution. We have
\begin{equation}\label{u:eq:scalar-bounds}
 m_0=0,\qquad |m(t,y)|<1,\qquad
 0\le\chi(t,y)\le1-m(t,y)^2\le1.
\end{equation}
The potential and $\chi$ are even in $y$, and $m$ and $b$ are odd.
These bounds follow from \cref{inp:ob-scalar} and the step-recursion
bound in \cite[Section~3]{oai-typical-upper}. For the latter bound, we note
that one Gaussian step at mass $\lambda\in[0,1]$ satisfies
$(Tf)'=\E_\lambda f'$ and
$(Tf)''=\E_\lambda f''+\lambda\Var_\lambda(f')$, where
$\E_\lambda$ is the normalized tilted expectation. If
$0<f''\le1-(f')^2$, then
$(Tf)''\le1-((Tf)')^2$.
By induction from $\log(2\cosh y)$ and scalar stability, we obtain
the inequality for $\chi$. Since $\chi>0$ by
\cref{inp:ob-scalar}(i), this implies that $|m|<1$. On the
continuation, the formulas are explicit.

\begin{proposition}[Scalar foundation]\label{u:prop:scalar}
Under the Wiener law in \eqref{u:eq:scalar-process}, $m_t$ is a bounded
martingale in the whole $Y$-past. Define
\begin{equation}\label{u:eq:scalar-profiles}
\begin{gathered}
 p(t)=\E\chi_t,\qquad r(t)=\E\chi_t^2,\qquad
 \varkappa(t)=\E\frac{\chi_t^2}{1-m_t^2},\\
 \rho(t)=\beta\sqrt{r(t)},\qquad
 \sigma(t)=\sqrt{1-\rho(t)^2}.
\end{gathered}
\end{equation}
All expectations are under the unlinked, zero-seed scalar Wiener law.
For $0\le s\le t$, we have
\begin{align}
 &\E(m_t-m_s)^2=\int_s^t r(v)\dd v,\qquad
   \int_0^t r(v)\dd v=\E m_t^2\le1,
   \label{u:eq:scalar-variance}\\
 &r=\beta^{-2},\quad\rho=1,\quad\sigma=0\quad(t\le v_*),
   \qquad 0<r(t)\le\beta^{-2},\quad p(t)^2\le r(t),
   \label{u:eq:scalar-stationarity}\\
 &r\text{ is nonincreasing},\quad |r'|\le2r,\qquad
   p\text{ is nonincreasing},\quad |p'|\le p
   \quad\text{almost everywhere}.
   \label{u:eq:profile-variation}
\end{align}
After $v_*$, we have $p=1-\E m_t^2$, $\varkappa=p$, and $p'=-r$.
The functions $\rho$ and $\sigma$ are continuous and of bounded
variation on every compact interval. Define
\begin{equation}\label{u:eq:scalar-energy-information}
 e(V)=\beta\int_0^V r(t)p(t)\dd t,\qquad
 j(V)=\frac12\int_0^V\rho(t)^2\varkappa(t)\dd t.
\end{equation}
For $V\ge v_*$,
\begin{equation}\label{u:eq:scalar-balance}
 j(V)-\log2+P_\beta-\beta e(V)=\frac{\beta^2}{4}p(V)^2,
 \qquad \E_{\gamma_n}F^c\le nP_\beta.
\end{equation}
For a fresh solution started at $Y_s=y$, we also have
\begin{equation}\label{u:eq:restart-coordinate}
 \E[Y_t m(t,Y_t)]-ym(s,y)
 =\int_s^t\E[\chi(v,Y_v)+m(v,Y_v)b(v,Y_v)]\dd v.
\end{equation}
This holds separately before and after $v_*$, and its integrand is
identically one after $v_*$.
\end{proposition}

\begin{proof}
We begin with the martingale and profile identities, which follow
from the scalar inputs. Then we identify the pressure normalization
in the entropy--energy identity and prove the restart formula.

By \cref{inp:ob-scalar}(iv), with $v=\beta^2t$, the process $m$ is a
bounded martingale before $v_*$. The explicit continuation has the
same property by It\^o's formula. The variance identity
\eqref{u:eq:scalar-variance} follows from the scalar evolution in
\cref{inp:ob-scalar}(iv) and the continuation. By
\cref{cor:inp-identity}, $\E m_v^2=v/\beta^2$ on $[0,v_*]$.
Continuity implies that $r(v)=\beta^{-2}$ on that closed interval.
Using Jensen's inequality and \eqref{u:eq:scalar-bounds}, we find
that $p^2\le r\le p$.

For $0\le v\le v_*$, the identity
\eqref{eq:inp-stationarity} and $\Gamma(v/\beta^2)=v/\beta^2$ imply
\begin{equation}\label{u:eq:mean-susceptibility}
 p(v)=1-\int\max(q,v/\beta^2)\,\mu_\beta(\dd q).
\end{equation}
Differentiating this Lipschitz integral almost everywhere, we obtain
$p'=-\bar\alpha/\beta^2=-\bar\alpha r$. This also follows directly
from \cref{inp:ob-scalar}(iv). After $v_*$, we have
$\chi_t=1-m_t^2$, so that $p=1-\E m_t^2$, $\varkappa=p$, and
$p'=-r$. In both intervals, $|p'|\le p$ because $r\le p$.
By the post-support estimate in \cite[Lemma~3]{oai-typical-upper},
$r$ is nonincreasing and $-2r\le r'\le0$; before
$v_*$ it is constant.
This proves \eqref{u:eq:scalar-stationarity} and
\eqref{u:eq:profile-variation}. Since $r$ is continuous and
monotone, the functions $\rho$ and $\sigma$ are continuous and have
bounded variation.

In variance time, \cref{inp:parisi,inp:unique} identify
\[
 P_\beta=U(0,0)-\frac1{2\beta^2}
                  \int_0^{\beta^2}v\bar\alpha(v)\dd v.
\]
This is the counting-measure normalization used in
\cite[Lemma~4]{oai-typical-upper}. The identity there reads
$\log2-j(V)+\beta e(V)=P_\beta-\beta^2p(V)^2/4$, which is
\eqref{u:eq:scalar-balance}; the endpoint $V=v_*$ follows by
continuity. The pressure inequality follows from \cref{inp:parisi}
and \cref{lem:model-completion}(c).

For the restart identity, the scalar evolution of
\cref{inp:ob-scalar}(v), after the time change, gives zero drift for
$m(v,Y_v)$. By the product rule \cite[Theorem~3.3.6]{karatzas-shreve-1991},
the drift of $Y_vm(v,Y_v)$ is $\chi(v,Y_v)+m(v,Y_v)b(v,Y_v)$. We use
localization and the bounded spatial derivatives of \cref{sec:scalar}
to justify the product rule under step approximation.
The stochastic integrand is $Y_v\chi(v,Y_v)+m(v,Y_v)$, which is
square integrable on every finite interval because $b$ is bounded.
The drift is bounded by two, so we may take expectations and pass
to the limit. This proves \eqref{u:eq:restart-coordinate}. On the
continuation, its
integrand is $\operatorname{sech}^2Y_v+\tanh^2Y_v=1$.
\end{proof}

\trilogyneedspace{12\baselineskip}
\begin{lemma}[Cross moments and an integral representation]\label{u:lem:cross}
For $0\le s\le t\le H$, a bounded measurable test $A$ of the
Brownian past through $s$, and deterministic $f\in L^2([s,t])$,
\begin{equation}\label{u:eq:cross-moment}
 \E\left[A m_H\int_s^t f(v)\dd M_v\right]
 =\int_s^t f(v)\E[A\chi_v]\dd v.
\end{equation}
Moreover, $m_t-m_s=\int_s^t\chi_v\dd M_v$, with the integral
constructed by predictable left sums in $L^2$ under the original
Wiener law. These assertions hold across $v_*$.
\end{lemma}

\begin{proof}
We begin with the tested identity for Brownian increments, which we
prove using the restart formula. We then pass to deterministic $L^2$ weights and
identify the integral by its second moment.

First let $A$ be a bounded Borel test of $Y_{[0,s]}$. The restart law
is the product of its $Y$-past law with a fresh Wiener driver,
followed by Picard restart. Multiplying
\eqref{u:eq:restart-coordinate} by $A$ and integrating against this
product law, we have
\[
 \E[A(Y_tm_t-Y_sm_s)]
 =\int_s^t\E[A(\chi_v+m_vb_v)]\dd v,
 \qquad b_v=b(v,Y_v).
\]
For unbounded endpoint observables, we clip the $y$ factor of
$ym(t,y)$ at $[-L,L]$ and take $L\to\infty$. The conditional first
moments are dominated by $|Y_s|+\sqrt{t-s}+t-s$, which is
integrable since $\E|Y_s|\le\sqrt s+s$. The drift
$\chi+mb$ is bounded by two, so Fubini's theorem applies to the signed
time integral. If the interval crosses $v_*$, we split there and
retain the original past test on the second interval. By continuity,
the split produces no endpoint term.

The martingale identity extends to integrable past tests by clipping,
with domination by twice their absolute value. In particular,
$\E[AY_s(m_t-m_s)]=0$, and for each $v\in[s,t]$,
$\E[Ab_vm_t]=\E[Ab_vm_v]$. Substituting
$M_t-M_s=Y_t-Y_s-\int_s^t b_v\dd v$, we obtain
\begin{equation}\label{u:eq:increment-cross-moment}
 \E[A m_t(M_t-M_s)]=\int_s^t\E[A\chi_v]\dd v.
\end{equation}
After completion, the sigma algebras of the Brownian past and of the
field past agree: the Picard solution is causal, and
$M_u=Y_u-\int_0^u b(v,Y_v)\dd v$ recovers the driver causally.
The recovery map into continuous paths is Borel by bounded joint
measurability and rational evaluations. Thus
\eqref{u:eq:increment-cross-moment} holds for Brownian-past tests.
Projecting $m_H$ to an intermediate time and using linearity, we
obtain \eqref{u:eq:cross-moment} for step functions $f$. By Brownian
isometry and Cauchy--Schwarz, we may pass to $L^2$ weights: the two
errors
are at most $\norm A_\infty\norm{f-f_\ell}_2$ and
$\norm A_\infty\sqrt{t-s}\norm{f-f_\ell}_2$.

For a refining deterministic partition $s=v_0<\dots<v_\ell=t$, let
$Z_\ell=\sum_i\chi_{v_i}(M_{v_{i+1}}-M_{v_i})$.
Predictable Gaussian isometry yields
$\E Z_\ell^2\to\int_s^t r(v)\dd v$.
Using \eqref{u:eq:increment-cross-moment} on each cell with past
test $\chi_{v_i}$, and projecting $m_t$ to that cell's endpoint,
we also have
$\E[(m_t-m_s)Z_\ell]\to\int_s^t r(v)\dd v$.
In these limits, we used that $\chi$ is continuous along the
paths and $\chi\le1$. By \eqref{u:eq:scalar-variance},
$\E|Z_\ell-(m_t-m_s)|^2\to0$, which proves the representation.
\end{proof}

On a pair $(M,\widetilde M)$ of independent Wiener processes, we define
\begin{equation}\label{u:eq:reflection}
 B=\int\rho\dd M+\int\sigma\dd\widetilde M,\qquad
 D=\int\sigma\dd M-\int\rho\dd\widetilde M.
\end{equation}
For a deterministic continuous function $f$ of bounded variation,
the integral is the path function
$I_f(b)(t)=f(t)b(t)-\int_{(0,t]}b(v)\dd f(v)$ on paths starting at
zero. By step approximation and Gaussian isometry, it agrees with the
$L^2$ integral. The reflection matrix
$\left(\begin{smallmatrix}\rho&\sigma\\\sigma&-\rho\end{smallmatrix}\right)$
is symmetric and orthogonal at each time. It follows that $(B,D)$
has the product Wiener law and has the causal inverse
\[
 M=\int\rho\dd B+\int\sigma\dd D,\qquad
 \widetilde M=\int\sigma\dd B-\int\rho\dd D.
\]
These identities hold for the laws of the full pairs of paths and for
the inverse maps.
For step weights, the reflection squares to the identity on every
independent Gaussian increment. Passing to uniform limits of paths,
with integration by parts for functions of bounded variation, we
find that the identities hold for $\rho$ and $\sigma$. Taking a countable
intersection over rational times, we identify the continuous versions.
The future increments of $(B,D)$ are independent of the pasts of both
original drivers, since the approximating deterministic integrals use
only fresh increments of $(M,\widetilde M)$.

Let $\mathcal W$ denote Wiener measure. For a terminal sign $x$,
we define the one-row law by
\begin{equation}\label{u:eq:row-law}
 \nu_V(\dd x,\dd M,\dd\widetilde M)
 =(1+x m_V)(\Unif_1\otimes\mathcal W\otimes\mathcal W)
                  (\dd x,\dd M,\dd\widetilde M).
\end{equation}
Averaging over the spin for each pair of paths, we see that this law is
normalized and that the law of the entire pair of paths is unchanged. By martingale projection, we conclude that
$\E_{\nu_V}[x\mid M_{[0,t]},\widetilde M_{[0,t]}]=m_t$ for
$t\le V$. We also use $\nu_V$ for its image on the history
$(x,Y,B,D)$ under the Picard solution map and \eqref{u:eq:reflection}.
For $0\le s\le t\le V$ and deterministic $f\in L^2([s,t])$,
by \cref{u:lem:cross} and the independence of the second original
driver, we have
\begin{equation}\label{u:eq:row-cross-moments}
\begin{aligned}
 \E_{\nu_V}\left[x\int_s^t f(v)\dd B_v\right]
 &=\int_s^t f(v)\rho(v)p(v)\dd v,\\
 \E_{\nu_V}\left[x\int_s^t f(v)\dd D_v\right]
 &=\int_s^t f(v)\sigma(v)p(v)\dd v.
\end{aligned}
\end{equation}
The tested conditional form of the first identity holds with any
bounded test $A$ of both pasts at $s$, with right side
$\int_s^t f(v)\rho(v)\E[A\chi_v]\dd v$.
For the $M$ term, we first integrate the independent
$\widetilde M$ past; for the $\widetilde M$ term, we first
integrate its centered fresh future. Unbounded Brownian tests are
handled by clipping in $L^2$. Under $\nu_V$, the path marginal remains
Wiener, but adjoining the terminal sign to the filtration need not
preserve the Brownian property. The filtrations used later are
recorded in \cref{u:sec:proposal-laws}.

\trilogyneedspace{12\baselineskip}
\begin{lemma}[Terminal and inverse-profile bounds]\label{u:lem:tails}
There exist constants depending only on $\beta$ such that, for all
$t\ge0$,
\begin{equation}\label{u:eq:scalar-tails}
\begin{gathered}
 p(t)\le C_\beta\frac{e^{-t/2}}{\sqrt{1+t}},\qquad
 r(t)\ge c_\beta\frac{e^{-t/2}}{\sqrt{1+t}},\\
 R_V=\left(\min_{0\le t\le V}r(t)\right)^{-1}
 \le C_\beta e^{V/2}\sqrt{1+V}.
\end{gathered}
\end{equation}
In particular, $p(V)^2\le C_\beta e^{-V}$.
\end{lemma}

\begin{proof}
By \cite[Lemma~3, proof]{oai-typical-upper}, the post-support
transition kernel is
\[
 K_t(y,z)=e^{-t/2}\varphi_t(z-y)\frac{\cosh z}{\cosh y},
\]
where $\varphi_t$ is the centered Gaussian density of variance $t$.
It integrates to one, and direct differentiation gives the generator
$\partial_{zz}/2+\tanh(z)\partial_z$. Since the drift is bounded and
Lipschitz, uniqueness of the solution of the Picard equation
identifies its law. Let $\lambda_*$ be
the endpoint law at $v_*$. For $k>1/2$, we have
\[
 \E\chi_{v_*+t}^k=e^{-t/2}\int\lambda_*(\dd y)\operatorname{sech}y
       \int\varphi_t(z-y)\operatorname{sech}^{2k-1}z\dd z.
\]
For $k=1$ and $t\ge1$, the bounds
$\varphi_t\le(2\pi t)^{-1/2}$ and
$\int\operatorname{sech}z\dd z=\pi$ imply the upper bound.
For $k=2$, let $C'_\beta<\infty$ be such that
$\lambda_*([-C'_\beta,C'_\beta])>0$, and restrict the inner integral
to $|z|\le1$. For $t\ge1$ and $|y|\le C'_\beta$, we have
\[
 \varphi_t(z-y)\ge
 \frac{e^{-(C'_\beta+1)^2/2}}{\sqrt{2\pi t}},
\]
which implies the lower bound. On $0\le t\le1$,
\eqref{u:eq:profile-variation} implies that
$r(v_*+t)\ge\beta^{-2}e^{-2t}$. The finite initial interval is
absorbed into the constants, which depend only on the fixed
temperature. The inverse bound follows from the monotonicity of $r$.
\end{proof}
